\SourceLeaf{450}%
\par\Needspace{9\baselineskip}\begin{center}\textbf{BEMERKUNGEN.}\end{center}
\noindent\hspace*{1.2em}Die auf S.\ 397 dieses Bandes wiedergegebenen Mitteilungen an \textsc{Olbers} vom 9.\ Oktober 1825 zeigen, dass \textsc{Gauss} zu dieser Zeit damit begonnen hat, die Untersuchungen zu Papier zu bringen, die den Gegenstand seiner \textit{Disquisitiones generales circa superficies curvas} von 1827 bilden. Das oben abgedruckte Fragment \textit{Neue allgemeine Untersuchungen über die krummen Flächen} unterscheidet sich von den \textit{Disquisitiones} im geringeren Umfang des Stoffes, in der Art der Behandlung und in der Anordnung der Sätze. Dort nimmt \textsc{Gauss} an, dass die rechtwinkligen Koordinaten $x,y,z$ eines Punktes der Fläche als Funktionen von zwei beliebigen unabhängigen Veränderlichen $p$ und $q$ dargestellt sind, während er hier die geodätischen Koordinaten $s$ und $\theta$ als neue Veränderliche wählt. Hier beweist er zuerst, dass sich die Summe der drei Winkel eines Dreiecks, das auf einer beliebigen krummen Fläche durch kürzeste Linien gebildet wird, von $180^{0}$ um den Flächeninhalt des Dreiecks unterscheidet, das ihm bei der Abbildung durch parallele Normalen auf der Hilfskugel mit dem Radius $1$ entspricht. Daraus leitet er durch einfache geometrische Betrachtungen das fundamentale Theorem her, dass »bei der Übertragung der Flächen durch Abwicklung das Krümmungsmaß an jeder Stelle unverändert bleibt«. Dort zeigt er dagegen zuerst in \S\ 12, dass sich das Krümmungsmaß allein durch die drei Größen $E,F,G$ und ihre Ableitungen nach $p$ und $q$ ausdrücken lässt, woraus der Satz über die Erhaltung des Krümmungsmaßes als Korollar folgt. Erst viel später, in \S\ 20, beweist er ganz unabhängig davon den Satz über die Winkelsumme eines geodätischen Dreiecks.
\begin{flushright}\textsc{Stäckel}.\end{flushright}
